Figure \ref{fig:jeff311-dnp-yield} shows the percentage of the total delayed neutron yield that comes from the six \glspl{DNP} with the largest yield contribution, where the yield from each \gls{DNP} is defined by Equation \eqref{eq:summation-tdny}.
There are only 136 \glspl{DNP} present in JEFF-3.1.1, and they have a total delayed neutron yield of 0.01477.
This is much smaller than the total delayed neutron yields from other datasets, and thus leads each \gls{DNP} to have a larger relative contribution to the delayed neutron yield.

\begin{figure}
    \centering
    \includegraphics[width=0.5\linewidth]{images/jeff311/dnp_yield.png}
    \caption{The fraction of the delayed neutron yield from particular \glspl{DNP} from JEFF-3.1.1 nuclear data. The five \glspl{DNP} with the largest delayed neutron yield contributions approximately 56\% of the total delayed neutrons per fission event. The remaining 131 \glspl{DNP} provide approximately 44\% of the total delayed neutron yield.}
    \label{fig:jeff311-dnp-yield}
\end{figure}


Table \ref{tab:jeff311-yld-hl} shows how the total delayed neutron yield $\nu_d$ and average half-life $\bar{\tau}$ compares between the value from the $I$ \glspl{DNP} and the six \gls{DNP} groups.
The differences between both are well within the uncertainty bounds, indicating the six groups are sufficiently capturing the general behavior of the delayed neutrons.
Although the total delayed neutron yield is smaller than other datasets, the average half-life is longer, indicating that the JEFF-3.1.1 data may be missing or under-predicting concentrations of short-lived \gls{DNP} in particular.

\begin{table}[]
    \centering
    \caption{JEFF-3.1.1 nuclear data delayed neutron yield and average half-life values.}
    \begin{tabular}{lll}
        \hline
        Parameter & Value & Units\\
        \hline
        $\nu_d(I)$ & $0.0148 \pm 0.0008$ &$ [-]$\\
        $\nu_d(K)$ & $0.0148 \pm 0.0007$ & $[-]$\\
        $\Delta \nu_d$ & $1.25 \pm 108.88$ & $[pcm]$\\
        \hline
        $\bar{\tau} (I)$ & $9.2 \pm 0.4$ & $[s]$\\
        $\bar{\tau} (K)$ & $9.2 \pm 0.4$ & $[s]$\\
        $\Delta \bar{\tau}$ & $0.0 \pm 0.6$ & $[s]$\\ % 0.015 +/- 0.584
        \hline
    \end{tabular}
    \label{tab:jeff311-yld-hl}
\end{table}

Table \ref{tab:jeff311-PCCi} shows the ten largest $PCC_i$ terms, which reveals which \glspl{DNP} have the greatest impact on the \gls{DNP} group parameters $\nu_{d,k}$ and $\tau_k$.
The top two \glspl{DNP}, $^{137}$I and $^{87}$Br, strongly affect groups two and one as previously discussed, which is why they have the largest \glspl{PCC}.
Section \ref{sec:iaea-jendl5} has a more detailed discussion on what effects increase \glspl{PCC}.


\begin{table}[]
    \centering
    \caption{JEFF-3.1.1 nuclear data ten largest \glspl{PCC}.}
    \begin{tabular}{ll}
        \hline
       Nuclide & $PCC_{i}$ \\
       \hline
       $^{137}$I & 3.619275 \\
       $^{87}$Br & 3.498093 \\
       $^{88}$Br & 3.323438 \\
       $^{138}$I & 2.993978 \\
       $^{96}$Rb & 2.727061 \\
       $^{89}$Br & 2.304603 \\
       $^{98}$Y & 2.055720 \\
       $^{94}$Rb & 1.959133 \\
       $^{91}$Br & 1.919183 \\
       $^{90}$Br & 1.886707 \\
        \hline
    \end{tabular}
    \label{tab:jeff311-PCCi}
\end{table}

Table \ref{tab:jeff311-Ui} shows the ten largest $U_i$ terms, which corresponds to which \glspl{DNP} have the largest uncertainties relative to their impact on the \gls{DNP} group parameters $\nu_{d,k}$ and $\tau_k$.
There are several \glspl{DNP} that appear in Table \ref{tab:jeff311-Ui} that also appear in Table \ref{tab:jeff311-PCCi}, indicated by the bold values.
This highlighting indicates \glspl{DNP} that strongly affect the \gls{DNP} group parameters, meaning that, although their relative uncertainties may not be as large as the other \glspl{DNP} in the table, the uncertainties they do have lead to large differences in the \gls{DNP} group parameters and calculated delayed neutron emission rates.

\begin{table}[]
    \centering
    \caption{JEFF-3.1.1 nuclear data ten largest uncertainty-weighted \glspl{PCC}. Bold \glspl{DNP} are present in Table \ref{tab:jeff311-PCCi} as well, indicating they have a strong effect on \gls{DNP} group parameters.}
    \begin{tabular}{ll}
        \hline
       Nuclide & $U_{i}$ \\
       \hline
       \textbf{$^{98}$Y} & \textbf{0.387086} \\
       $^{93}$Br & 0.237471 \\
       $^{147}$Xe & 0.214894 \\
       \textbf{$^{96}$Rb} & \textbf{0.176298} \\
       $^{138}$Te & 0.171012 \\
       $^{94}$Br & 0.152192 \\
       \textbf{$^{91}$Br} & \textbf{0.148399} \\
       $^{158}$Nd & 0.144933 \\
       \textbf{$^{137}$I} & \textbf{0.141580} \\
       $^{111}$Mo & 0.134707 \\
        \hline
    \end{tabular}
    \label{tab:jeff311-Ui}
\end{table}

Figure \ref{fig:jeff311-bar-uiv} gives a more detailed view on which \gls{DNP} values specifically lead to the large $U_i$ values.
This offers slightly more insight than Table \ref{tab:jeff311-Ui}, which sums over the various \gls{DNP} values.
For instance, these results show that $^{98}$Y has a large contribution of its uncertainty-weighted sensitivity come from emission probability uncertainty, with a smaller contribution from its half-life uncertainty.
It can also be seen that $^{137}$I is not in Figure \ref{fig:jeff311-bar-uiv} even though it is in Table \ref{tab:jeff311-Ui}, indicating that its uncertainty-weighted sensitivity comes from the sum of several smaller components rather than one large term.
Specifically, 54.5\% of its $U_i$ comes from emission probability uncertainty, 42.8\% comes from concentration uncertainty, and the remainder comes from half-life uncertainty.


\begin{figure}
    \centering
    \includegraphics[width=0.5\linewidth]{images/jeff311/pcc-bar.png}
    \caption{The ten largest uncertainty-weighted \glspl{PCC} $U_{i,v}$ sorted by \gls{DNP} and associated \gls{DNP} value from ENDF/B-VII.1 nuclear data.}
    \label{fig:jeff311-bar-uiv}
\end{figure}


Figure \ref{fig:pcc-charts-jeff311} shows the $PCC_i$ and $U_i$ plotted on the chart of the nuclides.
These plots do not limit to only the top ten largest values, but instead show how all of the values compare among each other.
The values of $U_i$ are all smaller than those of $PCC_i$, because all of the relative uncertainties are less than 1.
There are similar patterns as discussed in Section \ref{sec:iaea-jendl5}, with horizontal peaks in the $PCC_i$ for bromine, rubidium, and iodine.

\begin{figure}[!ht]
\centering
\subfloat[The $PCC_i$, or total \gls{PCC}, calculated for each \gls{DNP}.]{
  \includegraphics[width=0.45\linewidth]{images/jeff311/chart_PCC.png}
}
\hspace{0.05\linewidth}
\subfloat[The $U_i$, or total uncertainty-weighted \gls{PCC}, for each \gls{DNP}.]{
  \includegraphics[width=0.45\linewidth]{images/jeff311/chart_PCC_uncertainty.png}
}
\caption{Chart of nuclide plots showing the \glspl{PCC} and uncertainty-weighted \glspl{PCC} using JEFF-3.1.1 nuclear data.}
\label{fig:pcc-charts-jeff311}
\end{figure}

\FloatBarrier